\subsection{Contributions}

The core problem this thesis addresses is the absence of a principled, 
hardware-aware understanding of how algorithm selection and implementation 
choices jointly affect the sensitivity and computational performance of compact 
monopulse radar systems for UAV DAA. While the constituent pieces have each 
been studied in isolation, system designers currently lack the evidence needed 
to make informed trade-offs between detection sensitivity, processing latency, 
and throughput when deploying radar DAA on small UAV platforms with constrained 
compute and power budgets.

This thesis addresses this problem through a systems-level investigation of the 
software design space of a GPU-based radar tracking pipeline for a small 
monopulse radar. The challenges of poor angular resolution and low signal 
strength under SWaP constraints are addressed by comparing signal processing 
and tracking algorithms for their ability to detect faint targets from small 
UAV platforms. These algorithms are implemented on GPU hardware, with 
throughput, latency, memory movement costs, and synchronisation overhead 
measured through integration testing and profiling of the complete pipeline. 
The effect of different hardware contexts on optimal pipeline design is also 
investigated. The specific contributions are:

\begin{itemize}
    \item A comparison of candidate detection, angle-estimation, and tracking 
    algorithm combinations, assessed by detection sensitivity and track quality 
    in multi-target aerial scenarios, to inform algorithm selection.
    \item A comparison of alternative low-level implementations of the selected 
    algorithms, quantifying their effects on end-to-end throughput and latency 
    in highly parallelised GPU pipelines, to inform implementation choices.
    \item An assessment of how these algorithm and implementation conclusions 
    vary across embedded hardware configurations, providing hardware-aware 
    design recommendations.
\end{itemize}